\scanpage{53}
\mspage{49}
% archive sheet number: 52
et via cette identification, le support d'une partie \inserted{$A$} de $\Sigma$ s'identifie à son Sup\inserted{$\xi$}, son cosupport à $C(\xi)$.

\deleted{Ceci dit, à un atelier \unsure{spatial} $(\unsure{\mathcal A|\mathscr A},\trianglelefteq,\mathrel{|\unsure{o|\circ}|})$ est associé un \unsure{spatios|Spatios} $(\Sigma,\leq,C)$, avec en plus une}

\unsure{Étant|Etant} donné un Spatios $\Sigma$, on \deleted{définit} appelle \uline{figure} \uline{\unsure{du|de}} \uline{Spatios}
\[
\deleted{\mathrm{Fig}\,\unsure{\mathrm{ind}|\mathrm{ord}}(\Sigma)}
\]
\[
\deleted{\mathrm{Fig}\,\mathrm{spat}(\Sigma)=\{(\deleted{\unsure{A|\mathcal A}}\unsure{\Phi|\varPhi},\rho)\mid\Phi\in}
\]
\unsure{un|une} \deleted{partie} couple $(\Phi,\rho)$, où
\[
\Phi\subset\Sigma
\]
est une partie de $\Sigma$, et $\rho$ une relation sur $\Phi$, satisfaisant les conditions

\uline{\deleted{\unsure{fig|Fig}}Fig spat 1}\quad $\rho$ est une relation d'ordre sur $\Phi$

\uline{Fig spat 2}\quad a) $\forall X\in\Phi$, on a $X\ne0_\Sigma$, et

\hspace*{4em}b) $\forall X,Y\in\Phi,\ X\ne Y\Longrightarrow X\mathrel{|o|}Y$ (i.e. $\unsure{x|X}\leq CY$)

\uline{NB} On pourrait être plus exigeant, en s'inspirant de l'axiome des supports \deleted{\unsure{de}}pour un atelier. On ne va pas s'en préoccuper pour le moment.

On désigne par $\mathrm{Fig}\,\mathrm{spat}(\deleted{\unsure{\mathcal A|A}}\Sigma)$ l'ensemble des figures spatiales de $\Sigma$, et par
